\section{研究问题与证据层级}

本报告把事实分成四层。第一层是公司定期报告、交易所公告和合作方官方资料；第二层是可归档的原始券商PDF；第三层是保留券商、日期和预测字段的公开一致预期快照；第四层是媒体摘要、生态合作和无法追到原件的说法。只有第一、二层可以支撑核心经营描述；第三层仅作盈利预期交叉检查；第四层不进入EPS、目标价或客户链事实。

这种分层尤其适用于智能汽车电子。OEM客户名字、芯片合作、产品定点和订单新闻通常含有真实的技术或商业进展，但它们不自动给出车型、SOP、单车ASP、收入份额、项目毛利或回款。若这些字段缺失，正确结论是“尚待验证”，而不是用行业常识填补。

\begin{exhibitbox}[来源层级与可用范围]
\small
\begin{tabularx}{\textwidth}{L{2.2cm}L{2.6cm}Y Y}
\toprule
层级 & 本案例材料 & 可用结论 & 明确不能证明 \\
\midrule
T1 官方披露 & 2025年报、2026Q1、Qualcomm与NVIDIA官方材料 & 分部收入/毛利、客户/平台的已披露量产或合作状态、股本与现金流 & 客户收入份额、车型ASP、未署名订单金额 \\
T2 原始券商PDF & 开源、东吴、中银、辉立4份归档报告 & 评级、预测日期、公开收入/利润/EPS假设 & 未写出的目标价、方法或客户证据 \\
T3 可审计快照 & 东方财富结构化一致预期 & 券商身份、日期、评级与EPS分布 & 目标价共识、估值方法、原始报告全文 \\
T4 弱来源 & 摘要、转载、搜索线索 & 后续搜证方向 & 任何正向估值权重 \\
\bottomrule
\end{tabularx}
\sourcenote{来源登记、券商索引和耗竭日志。}
\end{exhibitbox}

资料治理不是形式要求，而是模型的权限系统。若没有客户/车型/ASP/利润率的T1或等效证据，估值器不能把该关系带入基准。若没有当前可审计的券商目标价和估值法，Street权重必须为零；有EPS预测不等于有可用目标价。

\section{财务与市场数据的可审计边界}

2025年审计年报是历史财务的核心基准。2026Q1已经披露但未经审计，使用时只陈述已报数字，绝不年化。市场价格为2026年7月23日盘中快照：CNY83.48，前收CNY82.38，日内高低CNY83.49/CNY80.87；它不是最终收盘价，也不能用于构造未取得的历史收益、资金流或拥挤度结论。

\begin{exhibitbox}[关键数据质量]
\small
\begin{tabularx}{\textwidth}{L{3cm}Y Y}
\toprule
字段 & 可用状态 & 研究处理 \\
\midrule
2025收入、归母、分部毛利、现金流、股本 & 经审计且交叉核验 & 可进入历史桥与估值分母 \\
2026Q1收入、归母、营运资本 & 已披露、未经审计 & 仅用于趋势与下一季门槛，不年化 \\
2026H1正式披露/预告 & 本轮未取得 & 不作为预测已实现证据 \\
盘中价格、总市值、TTM P/E、P/B & 中高质量盘中快照 & 明示时点并用于目标价上行计算 \\
历史价格区间、北向、严格自由流通、限售 & not disclosed & 不作技术面、资金面或供给面的伪精确判断 \\
\bottomrule
\end{tabularx}
\sourcenote{已核验财务与市场数据包。}
\end{exhibitbox}

财务与市场数据的不足不会被另一种数据“补平”。例如，盘中P/E不等于前瞻P/E；年化订单不等于合同负债；客户合作不等于按客户的收入；生产量增加不等于每单位利润提高。每种指标只回答它被披露能够回答的问题。

\section{主张审计与信息耗竭}

来源审计保留了十二项高重要性主张，覆盖分部收入、客户量产、订单口径、平台合作、毛利、现金、公开盈利预测和历史目标价。五项耗竭记录说明：近期原始报告缺目标价、快照缺目标/营收/方法、辉立CNY147为过时历史线索、公共页面未定位可用原件、券商催化剂不能代替客户侧证据。

因此，读者不应将“没有找到”理解成“没有发生”，也不应将“公司未披露”理解成“可按乐观值估计”。本报告给出每个缺口的升级路径：公司公告、IR、OEM项目公告、正式中报或授权数据库导出。披露出现之前，估值保持在已确认经营基础上。
